\section{Bluetooth}
\label{sec:soft_bluetooth}

La sortie filaire décrite aux deux sections précédentes n'est qu'une des deux implémentations du \gls{sink}. La seconde est la sortie Bluetooth, objet de cette section. Elle diffère de la première sur un point déterminant. Là où la chaîne filaire pilote un matériel que le firmware possède de bout en bout, la sortie Bluetooth confie les échantillons à une pile protocolaire autonome qui possède la radio, impose sa propre cadence et n'accepte qu'un unique format. Le sink Bluetooth est donc autant une adaptation à ces contraintes qu'une simple destination d'écriture.

\subsection{Rôle et contraintes de plateforme}

Le lecteur émet le flux audio vers un récepteur extérieur, enceinte ou casque, au moyen du profil \gls{a2dp}. Ce profil distingue deux rôles, la source qui émet le flux et le récepteur qui le restitue~\cite{bluetooth_sig_a2dp}. Le lecteur tient donc ici le rôle de source, même si cette sortie est un sink du point de vue du firmware, puisqu'elle reçoit les échantillons du pipeline.

Le transport audio du Bluetooth classique repose sur un seul codec obligatoire, le \gls{sbc}~\cite{bluetooth_sig_a2dp}. Ce codec fixe deux propriétés du flux. Il travaille sur des échantillons de 16~bits, ce qui plafonne la résolution de la sortie Bluetooth là où la chaîne filaire atteint 24~bits. La fréquence d'échantillonnage est ensuite négociée une seule fois à l'ouverture de la session et ne peut plus changer d'une piste à l'autre, contrairement au convertisseur filaire.

Le firmware fixe cette fréquence à 44,1~kHz. Fixer une fréquence unique évite d'embarquer un rééchantillonneur ou de renégocier la session à chaque piste, deux solutions coûteuses en charge processeur comme en réactivité. Le choix de 44,1~kHz parmi les fréquences admissibles tient à deux raisons de commodité. Tout récepteur SBC est tenu de la prendre en charge~\cite{bluetooth_sig_a2dp} et elle est la fréquence native de la plupart des fichiers musicaux courants. Une piste d'une autre fréquence est refusée explicitement plutôt que restituée à une cadence fausse, un message le signalant à l'écran. Le rééchantillonnage est renvoyé aux évolutions possibles.

\subsection{Découplage et format des échantillons}
La pile Bluetooth fixe elle-même son rythme. Elle vient chercher les échantillons quand elle en a besoin, les encode en SBC et gère leur émission radio~\cite{ESP32_IDF_guide}. La sortie filaire fonctionne à l'inverse, puisque la tâche audio y pousse elle-même les échantillons vers le contrôleur \gls{dma}, comme le montre la figure~\ref{Chemins_sortie_audio.drawio}. Un tampon de découplage fondé sur les flux FreeRTOS~\cite{freertos_kernel_doc} s'intercale donc entre le pipeline qui produit les échantillons et la pile qui les consomme, comme annoncé à la section~\ref{sec:archi}.

\fig[H, width=0.92\textwidth]{Les deux chemins de sortie audio. Le pipeline alimente l'une ou l'autre sortie selon le sink actif. Diagramme créé avec l'aide de Claude Opus 4.8.}{Chemins_sortie_audio.drawio}

Ce tampon est placé en \gls{psram} car aucun contrôleur DMA n'y accède, le processeur assurant seul la recopie de son contenu. Dimensionné à environ 21~kio, il représente de l'ordre de 120~ms d'avance à 44,1~kHz sur des échantillons stéréo de 16~bits, coussin qui absorbe le caractère irrégulier de la consommation radio.

Le comportement aux deux bornes du tampon prolonge le mécanisme décrit à la section~\ref{sec:transport}. En cas de saturation, le sink n'accepte qu'une partie du bloc présenté et le pipeline présente de nouveau le reliquat au tour suivant. Lorsque le tampon se vide au contraire plus vite qu'il ne se remplit, la pile reçoit du silence en lieu et place des échantillons manquants, ce qui préserve la continuité de la liaison sans jamais bloquer la radio.

La représentation interne sur 32~bits justifiée à gauche doit enfin être ramenée aux 16~bits qu'attend l'encodeur. Le \gls{sink} ne conserve que les 16~bits de poids fort de chaque échantillon, soit un simple décalage. Le contenu utile étant justifié à gauche, l'opération n'écarte que des bits de poids faible, nuls pour une source de 16~bits et correspondant au dernier octet d'une source de 24~bits. Omettre cette conversion doublerait la cadence apparente et corromprait le flux, l'encodeur lisant alors chaque échantillon de 32~bits comme deux échantillons de 16~bits.

\subsection{Volume et garde anti-saturation}
\label{sec:volume_bt}

Le volume de la sortie Bluetooth obéit au même principe de commande unique que la sortie filaire, le potentiomètre restant la seule source de contrôle. Sa position est traduite dans l'échelle de volume absolu du profil \gls{avrcp} et transmise au récepteur pour application de l'atténuation~\cite{ESP32_IDF_guide}. Le firmware ignore volontairement les changements de volume émanant du récepteur afin de préserver le potentiomètre comme référence exclusive et d'éviter tout conflit de commande.

Comme le volume est appliqué par le récepteur, un nouveau risque apparaît qui n'avais pas lieu dans la chaîne filaire. Un récepteur qui vient de se connecter garde son propre réglage de volume interne, parfois très élevé, tant qu'il n'a pas reçu celui du lecteur. Envoyer des échantillons avant ce moment pourrait produire un son brutalement fort. Une garde bloque donc toute émission tant que le récepteur n'a pas confirmé la commande de volume. Elle joue le même rôle que l'amplificateur de la chaîne filaire, celui d'une porte qui ne s'ouvre qu'une fois la chaîne complète.

Cette garde complique la bascule du filaire vers le Bluetooth en cours de lecture. Le volume ne peut être envoyé qu'une fois la sortie routée vers le Bluetooth, mais l'audio ne peut démarrer qu'une fois le volume confirmé. La bascule est donc coupée en deux étapes. Elle est amorcée d'abord, puis la tâche de maintenance la termine dès la confirmation reçue. La séquence complète figure en annexe~\ref{ann:bascule-bt}. La sortie précédente continue de jouer pendant cette attente.

\subsection{Appairage, persistance et robustesse}

Une recherche liste les appareils à portée et l'utilisateur choisit celui auquel il souhaite se connecter~\cite{ESP32_IDF_guide}. Les appareils déjà connus sont gardés d'une session à l'autre dans une courte liste en \gls{nvs}, du plus récent au plus ancien. Oublier un appareil efface son entrée et le lien de sécurité qui va avec.

Allumer et éteindre la radio sont des opérations longues et bloquantes, confiées à des tâches éphémères selon le modèle de la section~\ref{sec:archi}. Deux comportements assurent ensuite la robustesse du lien. Si le lien tombe en cours de lecture, il est fermé proprement et la lecture reprend sur la sortie filaire, ce qui évite qu'une déconnexion coupe le son. Si une tentative de connexion échoue, la pile peut rester bloquée dans un état intermédiaire dont seule une remise à zéro complète de la radio la sort. Le firmware déclenche cette remise à zéro avant de reproposer la connexion. Les incidents rencontrés sur cette sortie et leurs corrections figurent en annexe~\ref{ann:fiabilisation}.